% Введение
\section{Introduction}

\begin{frame}
    \frametitle{Relevance of the research topic}
      \begin{itemize}
        \item \textbf{The nature and scale of the phenomenon:} the capitalist economy reproduces the conditions of general overproduction crises with regular periodicity, while cyclical dynamics permeate all spheres of economic life, requiring consideration in the development of economic theories and models, as well as in the implementation of state economic policy;
        \item \textbf{Historical variability and contemporary practice:} the historical development of the stages of capitalism and individual components of the capitalist economy do not eliminate the objective foundations of the cycle, but rather modifies its nature and manifestations;
        \item \textbf{Lack of systematization and classification of theories:} the existing systematization of theoretical concepts of the cycle in the scientific literature is incomplete, and the system of criteria for classifying and evaluating theories is methodologically flawed;
        \item \textbf{Misconceptions and ideological bias:} most of the existing cycle theories are predominantly positivist in nature and based on a metaphysical understanding of capitalism, describing the cycle primarily mechanistically through the lens of simple quantitative changes, rather than in the context of the system's dialectical self-development mediated by its immanent contradictions.
      \end{itemize}
\end{frame}

\begin{frame}
  \frametitle{Purpose and objectives of the research}
  \begin{columns}[T,onlytextwidth]
    \begin{column}{0.48\textwidth}
      \textbf{Purpose of the research:}
      \begin{itemize}
        \item Theoretical and empirical justification of the industrial cycle as a key specific form of the dialectical development of the capitalist economy at its current specific historical stage.
      \end{itemize}
    \end{column}

    \begin{column}{0.48\textwidth}
      \textbf{Objectives of the research:}
      \begin{enumerate}
        \item Define the concept of a cycle, systematize approaches to its definition and research;
        \item Identify the key elements of the cycle theory and systematize theoretical concepts based on them;
        \item Examine the mechanism of the industrial cycle from the perspective of dialectical development;
        \item Describe the reproduction of fixed capital as a key factor of the industrial cycle;
        \item Conduct concrete historical analysis of the crises and the cycle in the US economy;
        \item Perform an empirical analysis of indicators and measures of the industrial cycle in the US economy.
      \end{enumerate}
    \end{column}
  \end{columns}
\end{frame}

\begin{frame}
  \frametitle{Scientific novelty of the research (1)}
  \begin{enumerate}
    \item The most comprehensive systematization of various theoretical concepts of the cycle in the economy has been carried out to date, using the author's classification system developed within the framework of this work, based on key aspects of analysis and fundamental debatable issues of the cycle theory;
    \item The main and secondary contradictions of the capitalist economy and its constituent elements underlying the industrial cycle have been highlighted and specified, and their exacerbation and resolution during the crisis of overproduction have been described in detail. The role of the industrial cycle as a specific form of the dialectical development of the capitalist economy has been substantiated, where its main contradictions remain unchanged, while secondary contradictions are partially and temporarily resolved through force during the crisis of overproduction, and then reproduced again at a higher stage of development of the capitalist economy;
\end{enumerate}
\end{frame}

\begin{frame}
  \frametitle{Scientific novelty of the research (2)}
  \begin{enumerate}
    \setcounter{enumi}{2}
    \item Based on the principles of Marxist political economy, the author has developed a model of periodic mass renewal of fixed capital for the case of simple reproduction, vividly illustrating the mechanism of periodization and synchronization of the moments of renewal of fixed capital between various enterprises in the capitalist economy;
    \item Using a specific historical and political-economic analysis, all economic crises in the US economy have been systematized and classified based on their dates and typology, and the periods and phases of the industrial cycle in the US economy from 1919 to the present have been identified. The dynamics of the organic composition of capital and the rate of profit have been calculated and analyzed based on current statistical data, in relation to the phases of the industrial cycle, empirically confirming the fundamental principles of Marxist theory regarding the long-term trend of increasing the organic composition of capital and decreasing the rate of profit.
  \end{enumerate}
\end{frame}

\begin{frame}
  \frametitle{Theoretical and practical significance of the research}
  \begin{columns}[T,onlytextwidth]
    \begin{column}{0.48\textwidth}
      \textbf{Theoretical significance:}
      \begin{itemize}
        \item Further development of the Marxist theory of crises, improvement of the classification of economic cycle concepts, research on modern forms of capitalist cyclicity, and a deeper understanding of the historical limits of state regulation of the capitalist economy.
      \end{itemize}
    \end{column}

    \begin{column}{0.48\textwidth}
      \textbf{Practical significance:}
      \begin{itemize}
        \item A tool for the comprehensive analysis of the processes of reproduction of social capital, application in research and educational work, expert assessment of modern economic processes.
      \end{itemize}
    \end{column}
  \end{columns}
\end{frame}

\begin{frame}
  \frametitle{Area of research (1)}
  The content of the research, as well as the main provisions and conclusions, correspond to the passport of the VAK\footnote{Higher Attestation Commission} specialty 5.2.1 \textquotedbl{}Economic Theory\textquotedbl{}:

  \vspace{10pt}

  \begin{tabularx}{\textwidth}{l>{\raggedright\arraybackslash\hyphenpenalty=10000\exhyphenpenalty=10000}X}
    \textbf{p. 2} & Clarification of the economic content of the cycle category, determination of its structure, internal contradictions, and place in the system of categories of social capital reproduction, as well as differentiation of industrial, business, investment, credit, and financial cycles; \\
    \textbf{p. 4} & Development and application of a methodology for comprehensive research of the industrial cycle based on the dialectical and historical-materialistic method, the principle of moving from the abstract to the concrete, and the combination of theoretical, historical, and statistical levels of analysis; \\
    \textbf{p. 5} & Study of the historical development of cycle theories, identification of the main stages of their formation, and analysis of changes in the understanding of the causes and mechanisms of economic crises; \\
    \textbf{p. 6} & Systematization of scientific schools and research programs explaining the nature of economic cycles, as well as a comparative analysis of their subject, methodological and theoretical foundations;
  \end{tabularx}
\end{frame}

\begin{frame}
  \frametitle{Area of research (2)}
  \begin{tabularx}{\textwidth}{l>{\raggedright\arraybackslash\hyphenpenalty=10000\exhyphenpenalty=10000}X}
    \textbf{p. 7} & Concrete historical study of the emergence and evolution of the industrial cycle, the periodization of crises in the capitalist economy, and the comparison of cyclical dynamics at different stages of capitalist development; \\
    \textbf{p. 9} & Study of macroeconomic patterns of cyclical movement of the capitalist economy, including the relationship between the dynamics of production, accumulation, investment, profit, employment, credit and aggregate demand; \\
    \textbf{p. 11} & Application of political-economic approach to the study of the industrial cycle, in which cyclical dynamics are derived from the internal contradictions of production, accumulation, and reproduction of social capital; \\
    \textbf{p. 16} & Development of theoretical approaches to the study of the industrial cycle as a form of economic fluctuations in the capitalist economy, revealing its internal causes, phase structure, mechanism of movement and historically changing forms of manifestation.
  \end{tabularx}
\end{frame}